\documentclass[12pt,a4paper,oneside]{article}
\usepackage[brazil]{babel}
\usepackage[utf8]{inputenc}
\usepackage{olymp}


\newlength{\exmpwidouf}
\exmpwidinf=10cm
\exmpwidouf=6cm


\setcounter{problem}{3}

\begin{document}

\begin{problem}{Hilfe!}{hilfe}{2}

Para auxiliar os estrangeiros que vêm estudar na UFV, a DRI – Diretoria de Relações Internacionais –  preparou alguns panfletos com o significado de várias palavras. Sua missão é facilitar a utilização desses panfletos colocando as palavras em ordem alfabética.

%\Notes
%
%Observações, se tiver.

\InputFile

A entrada começa com uma linha contendo um número inteiro $N$. A próxima linha contém uma palavra identificando o idioma. Em seguida existem $N$ linhas, cada uma com duas palavras: a primeira é uma palavra em português e a segunda a tradução no idioma indicado. Restrições: $1 \leq N \leq 35$, as palavras contém apenas letras minúsculas (sem espaços) e cada uma tem no máximo 20 letras.

\OutputFile

A saída deve começar com uma linha com o texto ``\texttt{portugues -> xxx}'', sendo \texttt{xxx} o idioma indicado na entrada. Em seguida, as $N$ palavras da entrada, em ordem alfabética, com suas respectivas traduções fornecidas na entrada, conforme formato dos exemplos.

% \Notes
% Preste atenção aos tipos das variáveis, pois $F(90) \approx 2^{61}$.

\Example

\begin{example}
\exmp{
7
alemao
socorro hilfe
sim ja
nao nein
bom gut
obrigado danke
entrada eingang
saida ausgang
}{
portugues -> alemao
bom - gut
entrada - eingang
nao - nein
obrigado - danke
saida - ausgang
sim - ja
socorro - hilfe
}%
\end{example}

\begin{example}
\exmp{
3
bosnio
sim yup
entrada ulaz
cerveja pivo
}{
portugues -> bosnio
cerveja - pivo
entrada - ulaz
sim - yup
}%
\end{example}



%Comentários sobre o exemplo, se necessário.

\end{problem}

\end{document}
